\documentclass[11pt]{article}

% Change "review" to "final" to generate the final (sometimes called camera-ready) version.
% Change to "preprint" to generate a non-anonymous version with page numbers.
% \usepackage[review]{acl}
\usepackage[final]{acl}

% Match the stock ACL review-mode layout: line numbers always in the
% right margin of their column (col 1 in the column gap, col 2 in the
% page's right margin). One of the user-added packages flips lineno's
% switch state and this forces it back.
\makeatletter
\@ifpackageloaded{lineno}{%
  % Override running-mode placement: column 1 in left margin,
  % column 2 in right margin. Decision uses \if@firstcolumn which
  % is correct except at column-break boundaries; ~7 lines at those
  % boundaries may land in the inter-column gap rather than the
  % outer margin. Accepted tradeoff: avoids the column-gap-only
  % layout that comes from \rightlinenumbers*.
  \AtBeginDocument{%
    \def\makeLineNumber{%
      \if@firstcolumn\makeLineNumberLeft\else\makeLineNumberRight\fi}%
    \let\makeLineNumberOdd\makeLineNumber
    \let\makeLineNumberEven\makeLineNumber
    \let\makeLineNumberRunning\makeLineNumber
  }%
}{}
\makeatother

% Standard package includes
\usepackage{times}
\usepackage{latexsym}

% For proper rendering and hyphenation of words containing Latin characters (including in bib files)
\usepackage[T1]{fontenc}
% For Vietnamese characters
% \usepackage[T5]{fontenc}
% See https://www.latex-project.org/help/documentation/encguide.pdf for other character sets

% This assumes your files are encoded as UTF8
\usepackage[utf8]{inputenc}

% This is not strictly necessary, and may be commented out,
% but it will improve the layout of the manuscript,
% and will typically save some space.
\usepackage{microtype}

% This is also not strictly necessary, and may be commented out.
% However, it will improve the aesthetics of text in
% the typewriter font.
\usepackage{inconsolata}

%Including images in your LaTeX document requires adding
%additional package(s)
\usepackage{graphicx}

% Custom
\usepackage{xcolor}
\newcommand{\red}[1]{\textcolor{red}{#1}}
\usepackage{minted}

\newcommand{\Lucy}[1]{[\textcolor{orange}{Lucy: {#1}}]}
\newcommand{\Sally}[1]{[\textcolor{red}{Sally: {#1}}]}

\usepackage{booktabs, multirow} % for borders and merged ranges
\usepackage{soul}% for underlines
\usepackage{xcolor,colortbl} % for cell colors
\usepackage{changepage,threeparttable} % for wide tables
\usepackage{amsmath}
\usepackage{float}

\usepackage{CJKutf8} 
\usepackage{ragged2e}
% If the title and author information does not fit in the area allocated, uncomment the following
%
%\setlength\titlebox{<dim>}
%
% and set <dim> to something 5cm or larger.

\title{MortarBench: Evaluating Mortgage Loan Origination Agents}

% Author information can be set in various styles:
% For several authors from the same institution:
% \author{Author 1 \and ... \and Author n \\
%         Address line \\ ... \\ Address line}
% if the names do not fit well on one line use
%         Author 1 \\ {\bf Author 2} \\ ... \\ {\bf Author n} \\
% For authors from different institutions:
% \author{Author 1 \\ Address line \\  ... \\ Address line
%         \And  ... \And
%         Author n \\ Address line \\ ... \\ Address line}
% To start a separate ``row'' of authors use \AND, as in
% \author{Author 1 \\ Address line \\  ... \\ Address line
%         \AND
%         Author 2 \\ Address line \\ ... \\ Address line \And
%         Author 3 \\ Address line \\ ... \\ Address line}

% \author{First Author \\
%   Affiliation / Address line 1 \\
%   Affiliation / Address line 2 \\
%   Affiliation / Address line 3 \\
%   \texttt{email@domain} \\\And
%   Second Author \\
%   Affiliation / Address line 1 \\
%   Affiliation / Address line 2 \\
%   Affiliation / Address line 3 \\
%   \texttt{email@domain} \\}

\author{
 \textbf{Matthew Toles\textsuperscript{1}},
 \textbf{Yunan Lu\textsuperscript{1}},
 \textbf{Manav Munjal\textsuperscript{1}},
 \textbf{Bojun Liu\textsuperscript{1}},
\\
 \textbf{Yuanhao Deng}\textsuperscript{2},
 \textbf{Stephanie Selig\textsuperscript{2}},
 \textbf{Derek Rindner\textsuperscript{2}},
 \textbf{Cheng Li\textsuperscript{2}},
 \textbf{Zhou Yu\textsuperscript{1}}
\\
\\
 \textsuperscript{1}Columbia University,
 \textsuperscript{2}Tidalwave
\\
 \small{
   \textbf{Correspondence:} \href{mailto:mt3639@columbia.edu}{mt3639@columbia.edu}
 }
}

\begin{document}
\maketitle
\begin{abstract}

Loan origination is the process by which a lender creates a new loan, from application and underwriting through approval and funding.
This process serves a critical role in evaluating the eligibility and level of risk posed by an applicant.
Recently, firms have begun using mortgage loan agents to augment human loan officers, despite a lack of any public benchmark.
To fill this gap, we present MortarBench, a loan origination agent benchmark.
MortarBench uses a financial data synthesis and mutation pipeline to generate examples with broad edge case coverage that match real-world distributions and questions.
We find that state-of-the-art LLMs perform poorly, with closed-source models achieving at most 77.1\% exact match accuracy.
We also discover systematic biases in LLM perception of foreignness and treatment of informal value transfer systems.
Noting weaknesses in baseline model reasoning and oversensitivity, we introduce CRIT, a confidence calibration framework.
Our method increases accuracy to 80.5\% while improving human control over sensitivity and risk management.
The high error rate and biases in baseline models suggest they should not yet be deployed to loan origination without a human in the loop.



\end{abstract}

\section{Introduction and Related Work}

The U.S. mortgage industry generates over \$1.7 trillion through loan origination annually, necessitating a complex web of risk assessment and regulatory compliance \cite{mba2024forecast}. 
Mistakes during these decisions carry severe legal and financial consequences, including loan defaults, liabilities, and fines for non-compliance.
Although human labor is costly and automation is appealing, fully replacing it introduces significant risk.

\begin{figure}
    \centering
    \includegraphics[width=1.0\linewidth]{figs/main.png}
    \caption{Example MortarBench task. A mortgage underwriter is provided with a bank statement and ULAD applicant profile. The underwriter asks a question pertaining to the documents to the chatbot assistant, which answers based on the same documents.}
    \label{fig:main}
\end{figure}

The loan origination process focuses on validating specific facts about an applicant’s financial status (Figure \ref{fig:main}) against their claims according to guidelines set by governing bodies such as Fannie Mae and Freddie Mac. 
Underwriters primarily rely on 1) standardized documents containing the borrower's self-reported profile (assets, liabilities, accounts, marital status, etc.), and 2) official bank statements. 
% The officer's goal is to ensure the structured claims in the ULAD—such as income, assets, and recurring debts—align perfectly with the ground truth found in the transaction records.
The officer's goal is to ensure consistency between the profile and bank statements, and that the application adheres to regulatory guidelines.
However, analyzing the application package, and especially bank statements, is a difficult and tedious process.
Because bank statements and transaction descriptions do not follow any standard format, reconciliation for a question such as ``What is the total value of the applicant's unsecured loans?'' is in fact a complex natural language and numerical task.

Recently, loan origination firms have begun integrating Large Language Model (LLM) agents into their workflows.
Despite this rapid integration, the industry lacks a standardized benchmark to evaluate these agents' performance.
This evaluative vacuum hinders broader adoption, stifles iterative development, and prevents objective comparisons between competing products.
% Ultimately, institutions are deploying "black-box" models into a highly regulated financial environment with no clear understanding of their true error rates or failure modes.
To address this gap, we propose MortarBench, a novel evaluation framework designed specifically for mortgage underwriting tasks.

Generating realistic, challenging examples of expert tasks in MortarBench presents two major challenges.
First, real financial data generally cannot be shared publicly due to privacy concerns, so we must generate internally consistent synthetic data that matches real world distributions.
Second, we must find a way to steer synthetic examples to cover rare edge cases relevant to user questions.
We overcome both challenges using a mutation-based synthetic data generation pipeline that allows us to procedurally edit input documents to force a desired answer while maintaining internal consistency.

Finding that baseline foundation models achieve weak accuracy on MortarBench, we contribute the CRIT Agent, confidence-calibrated model that improves accuracy, bias, and steerability.
CRIT improves up to 3.4\% over baseline foundation models (14.8\% error reduction).

% \section{Related Work}
Prior work has explored financial regulatory QA, including \citet{sohn2021global} and \citet{chen2024fintextqa}. 
Work on financial numeric and table QA includes \citet{zhu2021tat, chen2021finqa, reddy2024docfinqa, islam2023financebench}.
\citet{choi2025finder} analyzes RAG QA.
\citet{loukas2022finer} studies entity recognition in financial reports.
\citet{mollaev2025multimodal} releases large-scale anonymized transaction, geo-position, and technical support chat data from a major bank.
In contrast, we contribute a benchmark in the mortgage origination domain focused on the chatbot assistant role.
% Numerous works have explored the ability of LLMs to self-correct without external feedback, including \citet{han2024small, lu2023self, zhang2025ascot, sun2024enhancing, wang2022self, madaan2023self}.
Numerous works have explored the ability of LLMs to self-calibrate confidence in their own answers, including \citet{manakul2023selfcheckgpt, zhu2023calibration, kapoor2024calibration, kuhn2023semantic, kadavath2022language}.

We will release our code and data on GitHub upon acceptance.





% The check instruction is conditioned on the answer type to encourage precise verification: for boolean questions, the model must locate and quote the exact field or sentence that supports its Yes/No decision; for transaction-list questions, the model must quote the description, amount, and date of every candidate transaction to confirm it qualifies, remove any transaction that cannot be justified, and surface any qualifying transaction that was omitted;
% for account-list questions, the analogous check is applied at the account level.


\begin{figure}[h]
    \centering
    \includegraphics[width=0.5\textwidth]{figs/dataset.png}
    \caption{Dataset generation pipeline. We generate transactions across one or more bank statements to match real-world distributions, then create a user profile (ULAD) that conforms thereto. We filter and rewrite questions derived from users of an in-production chatbot assistant. For each question, we manually create a mutation function that edits the transaction and/or ULAD to conform to a desired answer.}
    \label{fig:pipeline}
\end{figure}

\section{MortarBench Creation Process}

To address the lack of standardized evaluation in the mortgage industry, we introduce MortarBench --- \textbf{Mor}tgage \textbf{T}ransaction, \textbf{A}nalysis, and \textbf{R}easoning Benchmark --- a specialized benchmark for evaluating LLM agents on realistic questions from loan officers.



Each data instance in MortarBench (Figure \ref{fig:pipeline}) consists of the following components:
% \begin{enumerate}
%     \item \textbf{Bank Statement:} A JSON-formatted list of financial transactions spanning 60 days.
%     \item \textbf{ULAD File:} A industry-wide standardized form containing the summary of applicants financial profile (e.g., employment history, stated assets, and liabilities).
%     \item \textbf{Target Question:} A single question about the bank statement and/or ULAD based on real-world loan officer workflows, curated and refined by expert annotators.
%     \item \textbf{Ground-Truth Answer:} The answer corresponding to the question, expressed as either a boolean value, a list of transaction identifiers, or a list of account identifiers.
% \end{enumerate}

\noindent \textbf{Bank Statement:} A JSON-formatted list of financial transactions spanning 60 days.

\noindent \textbf{ULAD File:} An industry-wide standardized form containing a summary of the applicant's financial profile (e.g., employment history, stated assets, and liabilities).

\noindent \textbf{Target Question:} A single question about the bank statement and, optionally, ULAD (standardized applicant financial summary), based on real-world loan officer workflows, curated and refined by expert annotators.

\noindent \textbf{Ground-Truth Answer:} The answer corresponding to the question, expressed as either a boolean value, a list of transaction identifiers, or a list of account identifiers.
% \subsection{Application Document Generation}




\subsection{Bank Statement Generation}
\label{sec:bank-statement}

To generate realistic transaction data that captures complex financial activity and supports diverse underwriting questions, we construct bank statements grounded in real financial records.
We sample 300 real bank statements from real loan applicants processed by a medium-sized loan processing company.
To overcome privacy limitations, we extract anonymous summary statistics of these applications, including the frequency and average value of key transaction types: payroll, benefits, and other income; deposits; housing and utilities expenses; bank fees; cryptocurrency transactions; and BNPL and other liabilities.
These statistics are used to generate synthetic bank statements with matching distributions.
Finally, we manually tag each transaction with transaction properties such as \texttt{unsecured loan} based on industry definitions.
Tags will be used later to calculate profile information and force specific answers in examples.
We limit bank statements to 60 days, aligning with regulatory guidelines.
We present an example bank statement in Appendix \ref{sec:app:bank-statement}.

\subsection{ULAD Generation}
\label{sec:ULAD-gen}

To generate ULAD documents consistent with the bank statement, we adopt a template-driven approach based on the standardized ULAD schema \citep{freddiemac_ulad}.
We populate each field relevant to dataset questions, including assets, collateral, liabilities, and loans with values computed from the bank statement using transaction property tags.
We present a summary of ULAD contents in Appendix \ref{sec:app:ulad}.

% To ensure consistency between financial activity and declared borrower information, we generate ULAD documents conditioned on the bank statement. We adopt a template-driven approach based on the standardized ULAD schema, which includes assets, collateral, liabilities, loan, party, and relationship fields. We populate these fields by directly mapping values from the bank statement to their corresponding structured entries. For example, asset accounts are instantiated from bank accounts, and liabilities are derived from repayment-related transactions. This design enforces cross-document consistency and produces application documents that reflect both borrower-reported information and verifiable financial activity.



% \subsection{Persona Creation}

% We construct nine personas from the Fannie Mae Test Credit Agency database \cite{FannieMae2017TestCreditReport}.
% We present the applicant data in the form of a standard Uniform Loan Application Dataset text file. 
% This file contains fields including the applicants assets, asset values, properties, liabilities, loan details, co-applicant information, demographics, and references.

% During the curation process, we minimally modify profiles to ensure a balanced distribution of positive and negative examples.
% This includes the intentional insertion of "rare-event" edge cases, such as discrepancies between the ULAD claims and bank statement reality. 

\subsection{Question Creation}
\label{sec:question-creation}

To generate realistic user questions, we collect real-world user questions made to an in-production loan assistant chatbot.
The chatbot is used to assist in QA tasks during the loan origination process.
We manually filter these questions to those that are 1) answerable using logical and financial reasoning or the Fannie Mae and Freddie Mac origination and underwriting selling guides \citep{fanniemae_selling_2025, freddie_mac_guide_2025} and 2) prevalent and useful within the loan origination process according to two subject matter expert authors. 
Subject matter expert authors then manually rewrite these questions to avoid multi-turn and temporal dependencies; ambiguous user intent; subjective questions; and questions that could leak personally identifiable information.

We find that answers fall into three categories:

\noindent \textbf{Boolean} - Answered with yes or no. Example: Do the payroll deposit entries match with the primary borrower's employer names stated in employment history?

\noindent \textbf{Transaction List} - Answered with a list of transaction IDs. Example: Which list of deposits, if any, are considered as large deposits?

\noindent \textbf{Account List} - Answered with a list of account IDs. Example: Which list of accounts, if any, are joint accounts where one or more account holders are not listed as borrowers on the loan application?

We also identify exactly one question requiring a monetary value answer.



\subsection{Application Files Mutation}

To generate balanced benchmark instances at scale, we design mutation functions that modify ULAD and bank statement files according to target questions and desired answers.
% We consider two types of mutations: single-file and multi-file mutations. Single-file mutations edit one document.
For example, for the question ``How many buy now, pay later (BNPL) transactions occur in the transaction list?'' with the desired answer ``3'', the mutation adds or removes existing BNPL transactions using the transaction tags defined in Section \ref{sec:bank-statement} until there are exactly three BNPL transactions.
% Multi-file mutations edit both the ULAD and bank statement when a question requires cross-document consistency checks, such as rental income consistency or undisclosed income detection.
If necessary, we update values in the ULAD, such as total liabilities, to reflect changes in the bank statement.
Unlike with fully random bank statement generation, this process creates a balanced dataset while preserving inter-document consistency.

\subsection{Summary Statistics}

We identify 47 unique questions satisfying the criteria in Section~\ref{sec:question-creation}.
For each question, we generate 4 cases, each with a profile consisting of a ULAD and a bank statement.
For boolean questions, we mutate half of the profiles to yield a ``yes'' answer and the other half to yield a ``no'' answer.
For transaction and account list questions, we mutate half of the profiles to contain a random non-zero number of transactions or accounts, and the other half to contain zero transactions or accounts.
This results in a balanced dataset containing 188 unique example scenarios.


\begin{table}[!htp]\centering\small
%\resizebox{\textwidth}{!}{ % use this if the table is too large
\begin{tabular}{lrr}\toprule
Total test cases &188 \\
Number of unique questions &47 \\
% Avg.\ bank statements per case &4.6 \\
Avg.\ transactions per bank statement &33.6 \\
ULAD fields &162 \\
% Choices per ULAD Field &7.4 \\
Questions needing ULAD &53.2\% \\
% Questions needing bank statement &100.0\% \\
% Questions needing both &53.2\% \\
Boolean answers &29.8\% \\
ID list answers &57.4\% \\
Account list answers &10.6\% \\
\bottomrule
\end{tabular}
\caption{Summary statistics for the MortarBench dataset.}\label{tab:mortarbench-stats}
\end{table}

\subsection{Evaluation Metrics}
We evaluate agent performance using a unified metric across question types.
To allocate partial credit to partially correct list questions, we employ the exact match as our primary metric, but also report F1 to capture agents' partially correct answers when required to list multiple transactions.
For boolean questions, we assign an F1 score of $1$ for correct predictions and $0$ otherwise.



 \section{Methods}

Noting weak baseline model performance, we introduce the \textbf{C}onfidence \textbf{R}eflection \textbf{I}nference for \textbf{T}ransactions (CRIT) Agent.
CRIT is motivated by the severe imbalance between false positives (32) vs. false negatives (0) in Gemini, indicating oversensitivity in the base model.
CRIT differs from the baseline implementation in that, for each transaction in its answer, the agent also generates a confidence score (1-5) indicating the probability the transaction is correctly included.
The confidence is based on the number and strength of assumptions that must be made to justify including the transaction.
Finally, we drop all transactions with confidence below a threshold $T=5$.
Confidence adjustment serves both to align the model away from oversensitivity and to improve steerability.
For example, a user may decrease the confidence threshold on a question where false negatives carry a high cost but predicted transactions are easy to verify.
We use Gemini 3.1 Pro for all CRIT implementations.
Full prompt templates for both the baseline and CRIT pipelines are provided in Appendix \ref{appendix:prompts}.

% CRIT Agent is designed to improve document grounding and reasoning over sets of transactions.
% After the agent produces an initial answer, the CRIT runs an additional verification round in which the original question, the initial answer, and the full source documents (bank statement and ULAD) are presented together with a targeted check instruction.
% To discourage hallucination, we instruct the agent to quote the description, amount, and date of every relevant candidate transaction to confirm qualification, and drop those that cannot be justified.

\section{Results}



Our results show that overall, Gemini 3.1 Pro performs strongest on MortarBench, both with and without the CRIT framework.
In the baseline case, it achieves 80.8\% F1 score, which improves to 83.6\% using CRIT (Table \ref{tab:eval_mt_combined_nostd}).


% Plain-text rendering of the table above:
%   Method            | F1   | EM  
%   ------------------+------+-----
%   Claude Sonnet 4.6 | 56.1 | 51.4
%   GPT-5.5           | 79.9 | 76.8
%   Gemini 3.1 Pro    | 80.8 | 77.1
%   CRIT              | 83.6 | 80.5
\begin{table}[H]
\centering
\small
\begin{tabular}{lrr}
\toprule
Method & F1 & EM \\
\midrule
Claude Sonnet 4.6 & 56.1 & 51.4 \\
GPT-5.5 & 79.9 & 76.8 \\
Gemini 3.1 Pro & 80.8 & 77.1 \\
CRIT & \textbf{83.6} & \textbf{80.5} \\
\bottomrule
\end{tabular}
\caption{Exact Match Accuracy (\%, mean over $N=3$ trials).}
\label{tab:eval_mt_combined_nostd}
\end{table}

As expected, improvement occurs entirely among transaction list questions (Table \ref{tab:eval_mt_by_question_type_nostd}), while results on boolean and account list questions remain identical.
We also see that Claude Sonnet 4.6 struggles especially with Boolean questions, achieving worse than random accuracy.



% Plain-text rendering of the table above:
%   Method            | Boolean | Txn IDs | Account IDs
%   ------------------+---------+---------+------------
%   Claude Sonnet 4.6 | 36.3    | 63.0    | 85.0       
%   GPT-5.5           | 94.6    | 72.3    | 93.3       
%   Gemini 3.1 Pro    | 97.0    | 72.2    | 96.7       
%   CRIT (ours)       | 97.0    | 77.0    | 96.7       
\begin{table}[H]
\centering
\small
\begin{tabular}{lrrr}
\toprule
Method & Boolean & Txn IDs & Account IDs \\
\midrule
Claude Sonnet 4.6 & 36.3 & 63.0 & 85.0 \\
GPT-5.5 & 94.6 & 72.3 & 93.3 \\
Gemini 3.1 Pro & \textbf{97.0} & 72.2 & \textbf{96.7} \\
CRIT (ours) & \textbf{97.0} & \textbf{77.0} & \textbf{96.7} \\
\bottomrule
\end{tabular}
\caption{Accuracy by Question Type (F1 \%, mean over $N=3$ trials).}
\label{tab:eval_mt_by_question_type_nostd}
\end{table}




\color{black}


% We analyze failure modes in Gemini 3 Pro on the naturally phrased questions to identify root causes of inaccuracies. 
% We observe that over half of errors (58.1\%) contain an extra transaction while only 29.8\% contain an extra transaction, suggesting a systematic bias for false positives.
% Overall, models tend to perform somewhat better on transaction list questions compared to boolean questions.
% One cause may be that boolean questions require reasoning over long range dependencies in the ULAD, whereas transaction list questions tend to be restricted to the bank statement only.






\section{Failure Analysis and Discussion}

Current frontier models show promise for improving accuracy and efficiency in the loan origination process in the near future.
Specifically, they can likely be deployed to assist with boolean and account list questions.
However, performance on transaction list questions is substantially weaker, indicating a continued need for human-in-the-loop oversight.

\subsection{Failure Analysis}

To identify root causes of failure, we manually annotate failure modes for baseline Gemini based on reasoning traces (Table \ref{tab:failure_modes_txn}).
The largest portion of errors (33.3\%) occur due to inaccurate transaction classification, for example, failing to recognize a BNPL transaction.
27.8\% of errors occur due to value matching, e.g., inaccurately comparing a sum of earnest money deposits to the stated amount in the ULAD.
Rather than making the arithmetic errors common in earlier models, Gemini often assumes that a value \textbf{exceeding} the stated ULAD value should be considered valid.
Domain knowledge errors (22.2\%) include falsely assuming that all wire transfers are international or that deposits by co-borrowers are automatically documented.
In 11.1\% of cases, Gemini misinterprets the constraints in the prompt, for example, classifying a one-time housing payment as recurring because it might recur in the future.

\begin{table}[H]
\centering
\small
\begin{tabular}{lrr}
\toprule
Category & Count & \% of Total \\
\midrule
Transaction Classification & 12 & 33.3 \\
Value Matching & 10 & 27.8 \\
Domain Knowledge & 8 & 22.2 \\
Constraint Misinterpretation & 4 & 11.1 \\
Unclear & 2 & 5.6 \\
\bottomrule
\end{tabular}
\caption{Failure-mode breakdown over $N=36$ wrong predictions on transaction-list questions for Gemini 3.1 Pro.}
\label{tab:failure_modes_txn}
\end{table}


% To determine whether errors occur due to false belief as opposed to 
Constructing our dataset with multiple variations of profiles for each question allows us to determine to what degree errors result stochastically from variation in profiles versus from false beliefs or other systematic errors.
For two versions of the same question, we define risk ratio as:

\begin{align*}
  \mathrm{RR} = \frac{\Pr(\text{2nd wrong}\mid\text{1st wrong})}{\Pr(\text{2nd wrong}\mid\text{1st right})} \approx 12.7
\end{align*}

\noindent for baseline Gemini.
We find that, given that the model answered incorrectly on the first version, the likelihood of an error on the second version is 12.7× higher than after a first-attempt success. 
This indicates that the overwhelming majority of errors are due to false beliefs in the model rather than stochasticity.

CRIT successfully aligns Gemini to reduce oversensitivity in transaction selection, reducing the false positive rate significantly at the expense of a negligible increase in false negatives (Table \ref{tab:add_vs_drop}).
In addition to improving overall EM and F1 compared to baseline, selective thresholding also improves steering and interpretability (Figure \ref{fig:fp_vs_fn}) --- the inclusion threshold can be chosen based on the real-world financial risk or human verification cost associated with each error.
Additionally, self-reflective confidence values can be used by users to prioritize human review.

\begin{figure}
    \centering
    \includegraphics[width=0.99\linewidth]{figs/fp_vs_fn.png}
    \caption{False negative and positive rate as a function of CRIT threshold $T$.}
    \label{fig:fp_vs_fn}
\end{figure}



% Plain-text rendering of the table above:
%   Method            | Qs   | FP/Q | Qs  | FN/Q
%   ------------------+------+------+-----+-----
%   Claude Sonnet 4.6 | 48.7 | 2.5  | 1.7 | 0.05
%   GPT-5.5           | 34.7 | 1.9  | 1.3 | 0.03
%   Gemini 3.1 Pro    | 36.7 | 1.1  | 0.0 | 0.00
%   CRIT (ours)       | 29.3 | 1.0  | 3.3 | 0.06
\begin{table}[H]
\centering
\small
\begin{tabular}{lrrrr}
\toprule
 & \multicolumn{2}{c}{Extras (FP)} & \multicolumn{2}{c}{Drops (FN)} \\
\cmidrule(lr){2-3} \cmidrule(lr){4-5}
Method & Qs & FP/Q & Qs & FN/Q \\
\midrule
Claude Sonnet 4.6 & 48.7 & 2.5 & 1.7 & 0.05 \\
GPT-5.5 & 34.7 & 1.9 & 1.3 & 0.03 \\
Gemini 3.1 Pro & 36.7 & 1.1 & 0.0 & 0.00 \\
CRIT (ours) & 29.3 & 1.0 & 3.3 & 0.06 \\
\bottomrule
\end{tabular}
\caption{Count of transaction list questions (Q) with at least one false positive (FP) or false negative (FN) transaction and the average number of FP or FN transactions per question.}
\label{tab:add_vs_drop}
\end{table}

\subsection{Bias}

\begin{figure*}
    \centering
    \includegraphics[width=1\linewidth]{figs/bias_foreign_plot.png}
    \caption{Rate at which wire transfers are classified as foreign as a function of language for company and personal names. * indicates that names have been transliterated into Latin script. $\Delta$ indicates the difference between the mean of all non-English languages in their native script and English.}
    \label{fig:name_bias}
\end{figure*}
For legal and ethical reasons, it is important that lending institutions do not discriminate against applicants based on protected classes.
Through manual review, we identify two questions in our data with elevated potential for biased decisions:

\noindent $\bullet$ Which list of deposits, if any, could be of foreign origin?

\noindent $\bullet$ Which list of deposits, if any, suggest funds from a private savings club or similar informal arrangement?

Additionally, we study whether US-based money transfer services such as PayPal are treated differently than non-US services across ten deposit- and withdrawal-related questions.

We probe bias by inserting transactions of the form 

\centering
ACH CREDIT - \{entity name\}\par
\justifying

into the bank statement and measuring whether the retrieval rates differ between the control and experimental groups.

\subsubsection{Bias Involving Non-English Names}
To test whether the LLM perceives non-English names as more foreign, we set \{entity name\} as either an English or a non-English personal or company name.
For personal names, we use the five most common masculine and feminine given names and the ten most common family names, paired arbitrarily, from Wikipedia's lists of most common given names and surnames \citep{wiki_given_names, wiki_surnames}.

We observe an extreme level of bias based on the language of origin of the name.
Across all three baseline models, wire transfers from English names are classified as foreign origin in only 13.3\% of cases (Figure \ref{fig:name_bias}).
Gemini and GPT-5.5 never classify these transactions as foreign origin.
However, non-English names induce a foreign origin classification in 77.0\% of cases. 
This effect is substantially more pronounced in non-Latin scripts (AR, HI, ZH: 93.3\%) compared to non-English Latin scripts (ES, FR: 52.5\%).
% Baseline models classify Arabic and Chinese names as foreign origin in 60/60 cases.
Gemini and Claude classify Arabic, Hindi, and Mandarin Chinese names as foreign origin in a staggering 60/60 cases.
We find that we can reduce the foreign classification bias on non-Latin scripts to the level of Spanish and French by transliterating into Latin script (AR\textsuperscript{*}, HI\textsuperscript{*}, ZH\textsuperscript{*}, 53.9\%).
Our agent, CRIT, shows less non-English bias than the corresponding Gemini 3.1 Pro baseline model in all cases, reducing the foreign origin classification gap between English and non-English personal names from 77.0\% to 54.0\%.

We perform a similar experiment where \{entity name\} is one of ten fictitious company names machine translated into target languages.
We observe even more bias overall, with an average baseline model foreign origin classification delta of 78.7\%.
Transliterating AR/HI/ZH company names appears less effective at reducing bias than when applied to personal names.
These results indicate that transactions involving non-English entities, and especially non-Latin scripts, are likely to receive unequal classification as foreign as compared to those with English names.
Like with personal names, applying our agent framework CRIT to Gemini reduces bias across all languages.

% \subsubsection{Bias Involving Ethnically-Associated Saving Organizations}

% We investigate whether models unequally identify private savings club deposits for ethnically-associated savings organizations.
% Numerous and diverse private financial organizations have emerged around the world to service community needs in saving, insurance, and money transfer, such as the Chinese \begin{CJK}{UTF8}{gbsn}会\end{CJK} (huì) or South Asian Chit systems.
% We identify 12 such systems as well as three non-ethnically linked systems listed in Appendix \ref{sec:app:private-savings}.
% We include them in wire transfer transaction template as above using Latin script, prompting models with the private savings club question.
% All models demonstrate perfect retrieval accuracy on both ethnically and non-ethnically linked savings organizations.
% However, when the savings organization is replaced with the name of a peer-to-peer money transfer system, which should \textbf{not} be returned, models produce relatively many false positives.
% US-based P2P systems including PayPal, Venmo, Western Union, and Zelle are erroneously classified as savings clubs 

\subsection{Bias Involving Non-US Peer-to-Peer Payments}

We study whether transactions through US-based electronic payment systems (EPSs) are treated differently than those through non-US-based systems.
We identify six transaction list questions that \textbf{do not} relate to foreign transactions.
For each question, we inject a transaction using the template above where \{entity name\} is PayPal, Venmo, Zelle, or Western Union, or one of eight primarily non-US systems (Appendix \ref{appendix:money-transfer-services}).
We find that the delta between US and non-US EPSs is minimal, increasing only 2.1\% across baseline models.
However, when non-US EPSs are replaced with the names of informal value transfer systems (IVTSs), such as hawala, or huikuan, the recall rate in Gemini increases by 27.1\% but in Claude decreases by 8.3\% (Figure \ref{fig:ivts).
These results indicate that users of IVTSs are likely to receive unequal treatment as compared to those using formal EPSs.
Our system, CRIT, shows the least difference between non-US EPSs and IVTSs.

\begin{figure}[H]
    \centering
    \includegraphics[width=1.0\linewidth]{figs/bias_money_transfer_categories.png}
    \caption{Recall rate of transactions involving US EPSs, non-US EPSs, and non-US IVTSs. $\Delta_1 = \text{non-US EPS} - \text{US EPS}$, $\Delta_2 = \text{IVTS} - \text{non-US EPS}$.}
    \label{fig:ivts}
\end{figure}


\section{Conclusion}

We introduce MortarBench, a realistic benchmark with broad edge case coverage for evaluating mortgage loan assistants.
We propose CRIT, an effective and human-centric strategy for reducing chronic oversensitivity in frontier models.
CRIT reduces oversensitivity as well as non-English naming and IVTS-related biases in baseline models.

% We do not see a significant correlation between the number of answers in the transaction and accuracy, suggesting ...?

\section{Limitations}

This dataset contains questions based on user interactions with current models.
As models advance, user behavior will likely evolve, altering the distribution of quesitons, answers, and biases.
This work addresses only the US mortgage origination process and will require adaptation to apply to use in other countries.
We choose the US origination process as a target dataset due to its economic size, standardization, and data accessibility.

% \section{Ethical Considerations}

% Accuracy and fairness in the mortgage loan origination process has a significant impact on applicants' access to capital and origination firms' risk profile.
% Models presented here 



% \red{
% To investigate whether errors are due to misidentification of the nature of transactions, we include ground truth tags in the prompt, simulating perfect information about transaction details, such as whether the transaction represents an undeclared liability.
% Adding these details, baseline Gemini exact match performance on transaction ID questions increases from 66.0\% to 70.4\%, indicating that 12.7\% of errors are due to hallucinating transaction details.
% }

\makeatletter
\ifacl@anonymize\else
\section*{Acknowledgements}
This research received funding from Tidalwave as well as DAPLab corporate support in the form of funding and/or compute from Amazon, IntellectAI, Infosys, Veris, Shopify, Microsoft, Thinking Machines, Dandy, Perplexity, and Daytona.
The views and conclusions presented here are those of the authors and should not be interpreted as representing the official positions of the funding organizations.
\fi
\makeatother

% Bibliography entries for the entire Anthology, followed by custom entries
%\bibliography{custom,anthology-overleaf-1,anthology-overleaf-2}

% Custom bibliography entries only
\bibliography{custom}

\appendix

\section{Example Bank Statement}
\label{sec:app:bank-statement}
\begin{minted}[fontsize=\small, breaklines]{json}
{
  "seed": "generated-test-15d7a496",
  "statement_type": "personal",
  "override_accounts": [
    {
      "type": "depository",
      "subtype": "savings",
      "starting_balance": 25701.58,
      "currency": "USD",
      "numbers": {
        "account": "54712641"
      },
      "transactions": [
        {
          "description": "Savings Club",
          "amount": 384.1,
          "currency": "USD",
          "transaction_id": "plaid-15d7a496-00024",
          "date_transacted": "2026-03-03",
          "date_posted": "2026-03-04"
        },
        {
          "description": "Loan Proceeds",
          "amount": 5671.38,
          "currency": "USD",
          "transaction_id": "plaid-15d7a496-00025",
          "date_transacted": "2026-03-13",
          "date_posted": "2026-03-14"
        }
      ],
      "identity": {
        "names": [
          "John Homeowner"
        ],
        "emails": [
          {
            "data": "john.homeowner@testmail.com",
            "primary": true,
            "type": "primary"
          }
        ],
        "addresses": [
          {
            "primary": true,
            "data": {
              "country": "US",
              "city": "Washington",
              "street": "175 13th St",
              "postal_code": "20013",
              "region": "DC"
            }
          }
        ]
      },
      "end_balance": 31757.06
    },
    {
      "type": "depository",
      "subtype": "checking",
      "starting_balance": 6153.25,
      "currency": "USD",
      "numbers": {
        "account": "82515638"
      },
      "transactions": [
        {
          "description": "Transfer from Alice Homeowner",
          "amount": 500,
          "currency": "USD",
          "transaction_id": "plaid-15d7a496-00022",
          "date_transacted": "2026-05-01",
          "date_posted": "2026-05-02"
        },
        {
          "description": "ACH DEBIT - Sezzle PMT",
          "amount": -65.74,
          "currency": "USD",
          "transaction_id": "plaid-15d7a496-01001",
          "date_transacted": "2026-04-27",
          "date_posted": "2026-04-28"
        },
        {
          "description": "Purchase at Crypto.com",
          "amount": -277.11,
          "currency": "USD",
          "transaction_id": "plaid-15d7a496-00007",
          "date_transacted": "2026-04-13",
          "date_posted": "2026-04-14"
        }
      // ...
      ]
    }
  // ...
  ]
}
\end{minted}

\section{ULAD Summary}
\label{sec:app:ulad}
A ULAD (Uniform Loan Application Dataset) is a standardized, machine-readable data format used in the U.S. mortgage industry to represent a residential loan application.
It includes borrower-level data (identity, demographics, income, employment, and legal declarations), financial position (assets such as bank accounts and liabilities such as outstanding debts), and property-related information (address, characteristics, valuation, and sales contract details).
It also captures the terms of the loan being originated, including loan amount, purpose, lien status, and interest features, along with metadata about the originating system.
Finally, ULAD represents relationships between entities (e.g., linking borrowers to assets and liabilities), enabling a complete, machine-readable representation of the loan application for automated processing and underwriting.

% \label{sec:appendix}

% This is an appendix.

\section{Money Transfer Services}
\label{appendix:money-transfer-services}

This appendix enumerates the money transfer services considered in our
analysis, grouped by category.

\begin{table}[H]
    \centering
    \begin{tabular}{lll}
        \toprule
        \textbf{US} & \textbf{Non-US} & \textbf{IVTS} \\
        \midrule
        Western Union & Alipay        & Hawala  \\
        Zelle         & M-Pesa        & Hundi   \\
        Venmo         & Orange Money  & Huikuan \\
        PayPal        & Paytm         & Padala  \\
                      & PIX           &         \\
                      & Wise          &         \\
                      & Revolut       &         \\
                      & Mercado Pago  &         \\
        \bottomrule
    \end{tabular}
    \caption{Money transfer services by category.}
    \label{tab:money-transfer-services}
\end{table}

\section{Prompts}
\label{appendix:prompts}

Table~\ref{tab:prompts} lists the prompt templates used by the baseline model
and by CRIT. Placeholders in braces (e.g.\ \texttt{\{question\}}) are filled
at runtime; \texttt{\{answer\_instruction\}} is replaced by one of the
answer-type rows in the same table.

\renewcommand{\arraystretch}{1.15}
\begin{table*}[h]
\centering
\small
\begin{tabular}{p{0.10\linewidth} p{0.13\linewidth} p{0.70\linewidth}}
\toprule
\textbf{Method} & \textbf{Stage} & \textbf{Prompt} \\
\midrule
Shared & Initial scaffold &
\texttt{\{question\}}\newline\newline
Bank Statement: \texttt{\{bank\_statement\}}\newline\newline
ULAD DU: \texttt{\{ulad\_du\}}\newline\newline
Answer the question. Do not think out loud. \texttt{\{answer\_instruction\}}. \\
\midrule
Baseline & Boolean (model) &
Answer with yes or no. \\
Baseline & Txn list (model) &
Describe the relevant transactions in text (titles/amounts/dates); do not guess or output transaction IDs. \\
Baseline & Account list (model) &
Identify the relevant accounts in text (names/descriptions/last4 digits); do not guess or output account IDs. \\
Baseline & Dollar (model) &
Think step by step about which dollar amounts are relevant to the question and why. Then identify and list each relevant amount from the documents, stating what it represents. Do not calculate totals yourself. \\
\midrule
Baseline & Boolean (clean) &
Question: \texttt{\{q\}}\newline
Unformatted answer: \texttt{\{raw\}}\newline
The answer given should be either yes or no. Read the question and answer, and simplify the answer to yes or no. Ignore any boilerplate; they are not part of the answer. \\
Baseline & Txn list (clean) &
Question: \texttt{\{q\}}\newline
Unformatted answer text (source of truth): \texttt{\{raw\}}\newline
Reference bank statement transactions JSON: \texttt{\{transactions\}}\newline
Step-by-step: (1) From the text only, count how many distinct transactions or payment occurrences are implied (N, allowing N+ if frequency suggests more). (2) Using the reference JSON, find all matching transactions; include additional matches if the pattern implies more than N. (3) If the text says none, return \texttt{[]}. Otherwise return ONLY a JSON list of all matching \texttt{TransactionID} values. Ignore boilerplate or conflicting IDs in the unformatted JSON. \\
Baseline & Account list (clean) &
Question: \texttt{\{q\}}\newline
Unformatted answer text (source of truth): \texttt{\{raw\}}\newline
Reference bank statement accounts JSON: \texttt{\{accounts\}}\newline
Use the text portion to decide which accounts the answer refers to. Match the mentioned account names/descriptions to the reference JSON and return ONLY a JSON list of the last 4 digits of the matching \texttt{AccountNumber} values. Ignore boilerplate or conflicting IDs. \\
Baseline & Dollar (clean) &
Question: \texttt{\{q\}}\newline
Unformatted answer: \texttt{\{raw\}}\newline
Return ONLY the final dollar amount as a plain number with two decimal places (e.g., 1234.56). No \$ sign, no commas, no other text. \\
\midrule
CRIT & Txn list (model) &
List every transaction in the bank statement that is plausibly related to the question, even if you are not sure. For each one, assign an integer confidence rating from 1 to 5: 1 = least likely to be relevant, 5 = clearly relevant. Do not mention transactions that are definitely irrelevant (confidence 0). Think out loud and state what assumptions you are making; a confidence of 5 should require no assumptions whatsoever. After stating your assumptions, return a JSON list of the form
\texttt{[\{"transaction\_id":\,"<id>","confidence":\,<1--5>\},\,\ldots]}, or \texttt{[]} if nothing is plausibly related. \\
CRIT & Txn list (clean) &
Question: \texttt{\{q\}}\newline
Unformatted answer text (source of truth): \texttt{\{raw\}}\newline
Reference bank statement transactions JSON: \texttt{\{transactions\}}\newline
The answer above should be a JSON list of
\texttt{\{"transaction\_id","confidence"\}} objects. Return ONLY a valid JSON list in EXACTLY that shape --- preserve every \texttt{transaction\_id} and its \texttt{confidence} value unchanged. If empty, return \texttt{[]}. \\
\bottomrule
\end{tabular}
\caption{Prompt templates for the baseline and CRIT methods. The ``model''
stage produces a free-form answer; the ``clean'' stage extracts a structured
final answer. CRIT differs from the baseline only in the transaction-list
prompts (model and clean); all other stages are shared. After the cleanup
pass, CRIT additionally filters the returned list in code, keeping only
transactions whose confidence is at least the threshold $T$.}
\label{tab:prompts}
\end{table*}

\end{document}
